\documentclass[11pt,letterpaper,oneside]{report}
\usepackage[T1]{fontenc}
\usepackage[utf8]{inputenc}
\usepackage{lmodern}
\usepackage[margin=0.90in,headheight=15pt]{geometry}
\usepackage{microtype}
\usepackage[svgnames,table]{xcolor}
\usepackage{amsmath,amssymb,booktabs,longtable,tabularx,array}
\usepackage{enumitem}
\usepackage[most]{tcolorbox}
\usepackage{listings}
\usepackage{fancyhdr}
\usepackage{titlesec}
\usepackage{hyperref}
\usepackage{xurl}
\usepackage{lastpage}
\definecolor{ink}{HTML}{14273C}
\definecolor{primary}{HTML}{0B5379}
\definecolor{accent}{HTML}{B66A30}
\definecolor{pale}{HTML}{F1F5F8}
\definecolor{codebg}{HTML}{F5F7F8}
\hypersetup{colorlinks=true,linkcolor=primary,urlcolor=primary,pdftitle={Linux and UNIX System Programming Engineering Suite}}
\titleformat{\chapter}[display]{\normalfont\LARGE\bfseries\color{ink}}{\color{accent}\large\MakeUppercase{\chaptername}\ \thechapter}{10pt}{}[\vspace{0.4ex}\titlerule]
\titleformat{\section}{\Large\bfseries\color{primary}}{\thesection}{0.7em}{}
\titleformat{\subsection}{\large\bfseries\color{ink}}{\thesubsection}{0.7em}{}
\setlist[itemize]{leftmargin=*,topsep=3pt,itemsep=2pt}
\setlist[enumerate]{leftmargin=*,topsep=3pt,itemsep=3pt}
\setlength{\parskip}{5pt}
\setlength{\parindent}{0pt}
\lstdefinestyle{sysc}{language=C,basicstyle=\ttfamily\footnotesize,backgroundcolor=\color{codebg},keywordstyle=\bfseries\color{primary},commentstyle=\color{ForestGreen},stringstyle=\color{Brown},breaklines=true,showstringspaces=false,tabsize=4,frame=single,rulecolor=\color{LightSteelBlue},columns=fullflexible,keepspaces=true}
\lstset{style=sysc}
\newtcolorbox{engineeringnote}[1][]{colback=pale,colframe=primary,boxrule=0.55pt,arc=1.5mm,title=Engineering note,fonttitle=\bfseries,#1}
\newtcolorbox{modernnote}[1][]{colback=white,colframe=accent,boxrule=0.55pt,arc=1.5mm,title=Modern extension --- not in the 2010 textbook,fonttitle=\bfseries,#1}
\newcommand{\sys}[1]{\texttt{#1}}
\pagestyle{fancy}\fancyhf{}\fancyhead[L]{\small\color{ink} Linux / UNIX Engineering Suite}\fancyhead[R]{\small\color{primary}\nouppercase{\leftmark}}\fancyfoot[C]{\small\thepage\ / \pageref{LastPage}}
\renewcommand{\headrulewidth}{0.45pt}
\setcounter{tocdepth}{1}
\setlength{\emergencystretch}{2em}
\begin{document}
\begin{titlepage}\centering\vspace*{2cm}{\color{primary}\Large ENGINEERING REFERENCE SERIES\par}\vspace{1.25cm}{\color{ink}\Huge\bfseries 64-Chapter Study Guides\par}\vspace{0.6cm}{\Large Original learning companions in the textbook's sequence\par}\vspace{2cm}\rule{0.75\textwidth}{1.2pt}\par\vspace{0.55cm}{\large An original companion to Michael Kerrisk\textquotesingle s 2010 text}\par\vspace{0.3cm}{\large Edition 1.0\quad |\quad October 2026}\par\vfill{\small Document 02\quad |\quad pdfLaTeX-compatible standalone source}\par\end{titlepage}
\tableofcontents\clearpage
\chapter{How to use these guides}
Follow the source chapters in order for a first pass, or follow the prerequisite graph in the curriculum. Each section below gives the source chapter title and printed opening page, key ideas to map while reading, an implementation focus, a failure mode to recognize, and an original verification task. These are independent study guides, not summaries or reproduced textbook exercises.

\section{Mandatory study loop}
\begin{enumerate}
\item Read the textbook chapter and identify the interface invariants and preconditions.
\item Map resource lifetimes and error return contracts.
\item Complete the challenge on an unprivileged Linux machine.
\item Write one failing or adversarial test, recording actual error codes.
\item Record whether the behavior is POSIX-specified, Linux-specific, or implementation-dependent.
\end{enumerate}
\chapter{Foundations and essential APIs}
\section{Chapter 1: History and Standards}
\textbf{Textbook entry:} printed page 1.\quad \textbf{Study label:} 2010 textbook foundation.

\subsection{Conceptual map}
\begin{itemize}
\item Unix lineage.
\item Abi versus api.
\item Standards versus implementation behavior.
\end{itemize}
\subsection{API and implementation audit}
Trace the purpose, argument lifetime, return-value convention and failure path for: uname, sysconf, getconf. Identify whether each name represents a kernel system call, a libc wrapper or a library-level operation on your platform; do not infer the category from naming alone.

\subsection{Failure mode and reasoning}
Assuming Linux behavior is portable to every UNIX. Construct a counterexample and explain which system contract or ownership invariant it violates.

\subsection{Practice and evidence}
Create an interface provenance table for ten calls and identify their governing standards. Record an expected result, a negative test, observed errno or pthread error code as applicable, and a cleanup trace.

\textbf{Review prompts.} Which state is process-local versus shared? What changes after an interruption, partial completion or concurrent operation?

\section{Chapter 2: Fundamental Concepts}
\textbf{Textbook entry:} printed page 21.\quad \textbf{Study label:} 2010 textbook foundation.

\subsection{Conceptual map}
\begin{itemize}
\item Kernel mediation.
\item Descriptor-centric i/o.
\item Processes and filesystem hierarchy.
\end{itemize}
\subsection{API and implementation audit}
Trace the purpose, argument lifetime, return-value convention and failure path for: open, read, fork. Identify whether each name represents a kernel system call, a libc wrapper or a library-level operation on your platform; do not infer the category from naming alone.

\subsection{Failure mode and reasoning}
Treating pathnames, inodes, and open file descriptions as interchangeable. Construct a counterexample and explain which system contract or ownership invariant it violates.

\subsection{Practice and evidence}
Draw a process/file-descriptor ownership diagram after a parent-child fork. Record an expected result, a negative test, observed errno or pthread error code as applicable, and a cleanup trace.

\textbf{Review prompts.} Which state is process-local versus shared? What changes after an interruption, partial completion or concurrent operation?

\section{Chapter 3: System Programming Concepts}
\textbf{Textbook entry:} printed page 43.\quad \textbf{Study label:} 2010 textbook foundation.

\subsection{Conceptual map}
\begin{itemize}
\item Return-value conventions.
\item Errno.
\item Libc wrappers.
\item Feature-test macros.
\end{itemize}
\subsection{API and implementation audit}
Trace the purpose, argument lifetime, return-value convention and failure path for: errno, perror, strerror\_r. Identify whether each name represents a kernel system call, a libc wrapper or a library-level operation on your platform; do not infer the category from naming alone.

\subsection{Failure mode and reasoning}
Failing to distinguish POSIX return codes from pthread error-number returns. Construct a counterexample and explain which system contract or ownership invariant it violates.

\subsection{Practice and evidence}
Build a tiny failure-reporting utility and test EACCES, ENOENT, and EINTR. Record an expected result, a negative test, observed errno or pthread error code as applicable, and a cleanup trace.

\textbf{Review prompts.} Which state is process-local versus shared? What changes after an interruption, partial completion or concurrent operation?

\section{Chapter 4: File I/O: The Universal I/O Model}
\textbf{Textbook entry:} printed page 69.\quad \textbf{Study label:} 2010 textbook foundation.

\subsection{Conceptual map}
\begin{itemize}
\item Open/read/write/close lifecycle.
\item File offsets.
\item Universality of descriptors.
\end{itemize}
\subsection{API and implementation audit}
Trace the purpose, argument lifetime, return-value convention and failure path for: open, read, write, close, lseek. Identify whether each name represents a kernel system call, a libc wrapper or a library-level operation on your platform; do not infer the category from naming alone.

\subsection{Failure mode and reasoning}
Assuming read or write transfers the entire requested count. Construct a counterexample and explain which system contract or ownership invariant it violates.

\subsection{Practice and evidence}
Implement a binary-safe copy loop with partial-write handling. Record an expected result, a negative test, observed errno or pthread error code as applicable, and a cleanup trace.

\textbf{Review prompts.} Which state is process-local versus shared? What changes after an interruption, partial completion or concurrent operation?

\section{Chapter 5: File I/O: Further Details}
\textbf{Textbook entry:} printed page 89.\quad \textbf{Study label:} 2010 textbook foundation.

\subsection{Conceptual map}
\begin{itemize}
\item Atomicity.
\item File description sharing.
\item Pread/pwrite.
\item Descriptor duplication.
\end{itemize}
\subsection{API and implementation audit}
Trace the purpose, argument lifetime, return-value convention and failure path for: fcntl, dup3, pread, pwrite. Identify whether each name represents a kernel system call, a libc wrapper or a library-level operation on your platform; do not infer the category from naming alone.

\subsection{Failure mode and reasoning}
Confusing per-descriptor flags with shared file-status flags. Construct a counterexample and explain which system contract or ownership invariant it violates.

\subsection{Practice and evidence}
Demonstrate shared offsets across dup and independent offsets across open. Record an expected result, a negative test, observed errno or pthread error code as applicable, and a cleanup trace.

\textbf{Review prompts.} Which state is process-local versus shared? What changes after an interruption, partial completion or concurrent operation?

\section{Chapter 6: Processes}
\textbf{Textbook entry:} printed page 113.\quad \textbf{Study label:} 2010 textbook foundation.

\subsection{Conceptual map}
\begin{itemize}
\item Pid identity.
\item Address-space layout.
\item Argv and environment.
\item Heap versus stack.
\end{itemize}
\subsection{API and implementation audit}
Trace the purpose, argument lifetime, return-value convention and failure path for: getpid, getenv, setjmp, longjmp. Identify whether each name represents a kernel system call, a libc wrapper or a library-level operation on your platform; do not infer the category from naming alone.

\subsection{Failure mode and reasoning}
Using pointer values across unrelated process address spaces. Construct a counterexample and explain which system contract or ownership invariant it violates.

\subsection{Practice and evidence}
Inspect memory regions and describe process environment inheritance. Record an expected result, a negative test, observed errno or pthread error code as applicable, and a cleanup trace.

\textbf{Review prompts.} Which state is process-local versus shared? What changes after an interruption, partial completion or concurrent operation?

\section{Chapter 7: Memory Allocation}
\textbf{Textbook entry:} printed page 139.\quad \textbf{Study label:} 2010 textbook foundation.

\subsection{Conceptual map}
\begin{itemize}
\item Heap allocation.
\item Allocation failure.
\item Lifetime and fragmentation.
\end{itemize}
\subsection{API and implementation audit}
Trace the purpose, argument lifetime, return-value convention and failure path for: malloc, calloc, realloc, free. Identify whether each name represents a kernel system call, a libc wrapper or a library-level operation on your platform; do not infer the category from naming alone.

\subsection{Failure mode and reasoning}
Integer overflow in allocation calculations and use-after-free. Construct a counterexample and explain which system contract or ownership invariant it violates.

\subsection{Practice and evidence}
Write an overflow-checked dynamic vector and test sanitizer findings. Record an expected result, a negative test, observed errno or pthread error code as applicable, and a cleanup trace.

\textbf{Review prompts.} Which state is process-local versus shared? What changes after an interruption, partial completion or concurrent operation?

\section{Chapter 8: Users and Groups}
\textbf{Textbook entry:} printed page 153.\quad \textbf{Study label:} 2010 textbook foundation.

\subsection{Conceptual map}
\begin{itemize}
\item Password and group databases.
\item Lookup semantics.
\item Reentrant variants.
\end{itemize}
\subsection{API and implementation audit}
Trace the purpose, argument lifetime, return-value convention and failure path for: getpwnam\_r, getgrgid\_r, getgroups. Identify whether each name represents a kernel system call, a libc wrapper or a library-level operation on your platform; do not infer the category from naming alone.

\subsection{Failure mode and reasoning}
Assuming identity lookup is fast, local, or reentrant. Construct a counterexample and explain which system contract or ownership invariant it violates.

\subsection{Practice and evidence}
Create a UID/GID inspection utility with dynamic buffers. Record an expected result, a negative test, observed errno or pthread error code as applicable, and a cleanup trace.

\textbf{Review prompts.} Which state is process-local versus shared? What changes after an interruption, partial completion or concurrent operation?

\section{Chapter 9: Process Credentials}
\textbf{Textbook entry:} printed page 167.\quad \textbf{Study label:} 2010 textbook foundation.

\subsection{Conceptual map}
\begin{itemize}
\item Real/effective/saved identities.
\item Privilege transitions.
\item Group memberships.
\end{itemize}
\subsection{API and implementation audit}
Trace the purpose, argument lifetime, return-value convention and failure path for: geteuid, setresuid, setgroups. Identify whether each name represents a kernel system call, a libc wrapper or a library-level operation on your platform; do not infer the category from naming alone.

\subsection{Failure mode and reasoning}
Mistaking effective-ID checks for a complete authorization model. Construct a counterexample and explain which system contract or ownership invariant it violates.

\subsection{Practice and evidence}
Trace the credentials of a controlled test process before and after a UID drop. Record an expected result, a negative test, observed errno or pthread error code as applicable, and a cleanup trace.

\textbf{Review prompts.} Which state is process-local versus shared? What changes after an interruption, partial completion or concurrent operation?

\section{Chapter 10: Time}
\textbf{Textbook entry:} printed page 185.\quad \textbf{Study label:} 2010 textbook foundation.

\subsection{Conceptual map}
\begin{itemize}
\item Wall versus monotonic clocks.
\item Calendar conversion.
\item Locale/timezone.
\end{itemize}
\subsection{API and implementation audit}
Trace the purpose, argument lifetime, return-value convention and failure path for: clock\_gettime, gmtime\_r, localtime\_r. Identify whether each name represents a kernel system call, a libc wrapper or a library-level operation on your platform; do not infer the category from naming alone.

\subsection{Failure mode and reasoning}
Scheduling intervals using a clock that can jump. Construct a counterexample and explain which system contract or ownership invariant it violates.

\subsection{Practice and evidence}
Build a monotonic elapsed-time measuring wrapper. Record an expected result, a negative test, observed errno or pthread error code as applicable, and a cleanup trace.

\textbf{Review prompts.} Which state is process-local versus shared? What changes after an interruption, partial completion or concurrent operation?

\section{Chapter 11: System Limits and Options}
\textbf{Textbook entry:} printed page 211.\quad \textbf{Study label:} 2010 textbook foundation.

\subsection{Conceptual map}
\begin{itemize}
\item Compile-time versus runtime bounds.
\item Indeterminate limits.
\end{itemize}
\subsection{API and implementation audit}
Trace the purpose, argument lifetime, return-value convention and failure path for: sysconf, pathconf, fpathconf. Identify whether each name represents a kernel system call, a libc wrapper or a library-level operation on your platform; do not infer the category from naming alone.

\subsection{Failure mode and reasoning}
Mistaking -1 with unchanged errno for an API failure. Construct a counterexample and explain which system contract or ownership invariant it violates.

\subsection{Practice and evidence}
Write a configurable buffer-sizing routine based on runtime constraints. Record an expected result, a negative test, observed errno or pthread error code as applicable, and a cleanup trace.

\textbf{Review prompts.} Which state is process-local versus shared? What changes after an interruption, partial completion or concurrent operation?

\section{Chapter 12: System and Process Information}
\textbf{Textbook entry:} printed page 223.\quad \textbf{Study label:} 2010 textbook foundation.

\subsection{Conceptual map}
\begin{itemize}
\item Process metadata.
\item /proc.
\item Uname and observation races.
\end{itemize}
\subsection{API and implementation audit}
Trace the purpose, argument lifetime, return-value convention and failure path for: uname, readlink, opendir. Identify whether each name represents a kernel system call, a libc wrapper or a library-level operation on your platform; do not infer the category from naming alone.

\subsection{Failure mode and reasoning}
Treating procfs snapshots as globally atomic. Construct a counterexample and explain which system contract or ownership invariant it violates.

\subsection{Practice and evidence}
Produce a bounded /proc status inspection tool. Record an expected result, a negative test, observed errno or pthread error code as applicable, and a cleanup trace.

\textbf{Review prompts.} Which state is process-local versus shared? What changes after an interruption, partial completion or concurrent operation?

\chapter{Filesystems, signals and timing}
\section{Chapter 13: File I/O Buffering}
\textbf{Textbook entry:} printed page 233.\quad \textbf{Study label:} 2010 textbook foundation.

\subsection{Conceptual map}
\begin{itemize}
\item User-space stdio.
\item Page cache.
\item Fsync and ordering.
\end{itemize}
\subsection{API and implementation audit}
Trace the purpose, argument lifetime, return-value convention and failure path for: setvbuf, fsync, fdatasync. Identify whether each name represents a kernel system call, a libc wrapper or a library-level operation on your platform; do not infer the category from naming alone.

\subsection{Failure mode and reasoning}
Assuming successful write means durable storage. Construct a counterexample and explain which system contract or ownership invariant it violates.

\subsection{Practice and evidence}
Compare buffered writes and sync behavior with explicit caveats. Record an expected result, a negative test, observed errno or pthread error code as applicable, and a cleanup trace.

\textbf{Review prompts.} Which state is process-local versus shared? What changes after an interruption, partial completion or concurrent operation?

\section{Chapter 14: File Systems}
\textbf{Textbook entry:} printed page 251.\quad \textbf{Study label:} 2010 textbook foundation.

\subsection{Conceptual map}
\begin{itemize}
\item Inodes.
\item Mount hierarchy.
\item Vfs.
\item Bind mounts.
\item Tmpfs.
\end{itemize}
\subsection{API and implementation audit}
Trace the purpose, argument lifetime, return-value convention and failure path for: statvfs, mount, umount2. Identify whether each name represents a kernel system call, a libc wrapper or a library-level operation on your platform; do not infer the category from naming alone.

\subsection{Failure mode and reasoning}
Assuming directory path identifies a stable mounted filesystem. Construct a counterexample and explain which system contract or ownership invariant it violates.

\subsection{Practice and evidence}
Read mountinfo and map mount points without modifying mounts. Record an expected result, a negative test, observed errno or pthread error code as applicable, and a cleanup trace.

\textbf{Review prompts.} Which state is process-local versus shared? What changes after an interruption, partial completion or concurrent operation?

\section{Chapter 15: File Attributes}
\textbf{Textbook entry:} printed page 279.\quad \textbf{Study label:} 2010 textbook foundation.

\subsection{Conceptual map}
\begin{itemize}
\item Ownership.
\item Modes.
\item Umask.
\item Stat.
\item Timestamps.
\end{itemize}
\subsection{API and implementation audit}
Trace the purpose, argument lifetime, return-value convention and failure path for: fstat, fchmod, fchown, utimensat. Identify whether each name represents a kernel system call, a libc wrapper or a library-level operation on your platform; do not infer the category from naming alone.

\subsection{Failure mode and reasoning}
Using access() then open() as an authorization check. Construct a counterexample and explain which system contract or ownership invariant it violates.

\subsection{Practice and evidence}
Build an fd-based metadata viewer and explain permission-check timing. Record an expected result, a negative test, observed errno or pthread error code as applicable, and a cleanup trace.

\textbf{Review prompts.} Which state is process-local versus shared? What changes after an interruption, partial completion or concurrent operation?

\section{Chapter 16: Extended Attributes}
\textbf{Textbook entry:} printed page 311.\quad \textbf{Study label:} 2010 textbook foundation.

\subsection{Conceptual map}
\begin{itemize}
\item Namespaced extended attributes.
\item Access policy and storage constraints.
\end{itemize}
\subsection{API and implementation audit}
Trace the purpose, argument lifetime, return-value convention and failure path for: listxattr, getxattr, setxattr. Identify whether each name represents a kernel system call, a libc wrapper or a library-level operation on your platform; do not infer the category from naming alone.

\subsection{Failure mode and reasoning}
Assuming xattrs are portable or always supported. Construct a counterexample and explain which system contract or ownership invariant it violates.

\subsection{Practice and evidence}
Inspect xattr support and handle ENOTSUP gracefully. Record an expected result, a negative test, observed errno or pthread error code as applicable, and a cleanup trace.

\textbf{Review prompts.} Which state is process-local versus shared? What changes after an interruption, partial completion or concurrent operation?

\section{Chapter 17: Access Control Lists}
\textbf{Textbook entry:} printed page 319.\quad \textbf{Study label:} 2010 textbook foundation.

\subsection{Conceptual map}
\begin{itemize}
\item Posix acl entries.
\item Mask.
\item Default acls.
\end{itemize}
\subsection{API and implementation audit}
Trace the purpose, argument lifetime, return-value convention and failure path for: acl\_get\_file, acl\_calc\_mask, acl\_set\_file. Identify whether each name represents a kernel system call, a libc wrapper or a library-level operation on your platform; do not infer the category from naming alone.

\subsection{Failure mode and reasoning}
Ignoring ACL masks when reasoning about effective permissions. Construct a counterexample and explain which system contract or ownership invariant it violates.

\subsection{Practice and evidence}
Construct a controlled ACL test fixture and verify access via getfacl. Record an expected result, a negative test, observed errno or pthread error code as applicable, and a cleanup trace.

\textbf{Review prompts.} Which state is process-local versus shared? What changes after an interruption, partial completion or concurrent operation?

\section{Chapter 18: Directories and Links}
\textbf{Textbook entry:} printed page 339.\quad \textbf{Study label:} 2010 textbook foundation.

\subsection{Conceptual map}
\begin{itemize}
\item Hard links.
\item Symlinks.
\item Rename.
\item Directory-fd resolution.
\end{itemize}
\subsection{API and implementation audit}
Trace the purpose, argument lifetime, return-value convention and failure path for: openat, unlinkat, renameat, readlinkat. Identify whether each name represents a kernel system call, a libc wrapper or a library-level operation on your platform; do not infer the category from naming alone.

\subsection{Failure mode and reasoning}
Following an untrusted symlink during pathname traversal. Construct a counterexample and explain which system contract or ownership invariant it violates.

\subsection{Practice and evidence}
Implement a directory-fd relative file reader with no-follow behavior. Record an expected result, a negative test, observed errno or pthread error code as applicable, and a cleanup trace.

\textbf{Review prompts.} Which state is process-local versus shared? What changes after an interruption, partial completion or concurrent operation?

\section{Chapter 19: Monitoring File Events}
\textbf{Textbook entry:} printed page 375.\quad \textbf{Study label:} 2010 textbook foundation.

\subsection{Conceptual map}
\begin{itemize}
\item Directory change notification.
\item Inotify event queue behavior.
\end{itemize}
\subsection{API and implementation audit}
Trace the purpose, argument lifetime, return-value convention and failure path for: inotify\_init1, inotify\_add\_watch, read. Identify whether each name represents a kernel system call, a libc wrapper or a library-level operation on your platform; do not infer the category from naming alone.

\subsection{Failure mode and reasoning}
Assuming every filesystem event is delivered without queue overflows. Construct a counterexample and explain which system contract or ownership invariant it violates.

\subsection{Practice and evidence}
Capture events and deliberately detect IN\_Q\_OVERFLOW handling. Record an expected result, a negative test, observed errno or pthread error code as applicable, and a cleanup trace.

\textbf{Review prompts.} Which state is process-local versus shared? What changes after an interruption, partial completion or concurrent operation?

\section{Chapter 20: Signals: Fundamental Concepts}
\textbf{Textbook entry:} printed page 387.\quad \textbf{Study label:} 2010 textbook foundation.

\subsection{Conceptual map}
\begin{itemize}
\item Signal dispositions.
\item Masks.
\item Pending signals.
\item Sigaction.
\end{itemize}
\subsection{API and implementation audit}
Trace the purpose, argument lifetime, return-value convention and failure path for: sigaction, sigprocmask, kill. Identify whether each name represents a kernel system call, a libc wrapper or a library-level operation on your platform; do not infer the category from naming alone.

\subsection{Failure mode and reasoning}
Treating signals as reliably queued state changes. Construct a counterexample and explain which system contract or ownership invariant it violates.

\subsection{Practice and evidence}
Build a minimal termination handler that only sets a flag. Record an expected result, a negative test, observed errno or pthread error code as applicable, and a cleanup trace.

\textbf{Review prompts.} Which state is process-local versus shared? What changes after an interruption, partial completion or concurrent operation?

\section{Chapter 21: Signals: Signal Handlers}
\textbf{Textbook entry:} printed page 421.\quad \textbf{Study label:} 2010 textbook foundation.

\subsection{Conceptual map}
\begin{itemize}
\item Async-signal-safety.
\item Sa\_siginfo.
\item Interrupted system calls.
\end{itemize}
\subsection{API and implementation audit}
Trace the purpose, argument lifetime, return-value convention and failure path for: sigaction, sigaltstack, write. Identify whether each name represents a kernel system call, a libc wrapper or a library-level operation on your platform; do not infer the category from naming alone.

\subsection{Failure mode and reasoning}
Calling printf or malloc from an asynchronous signal handler. Construct a counterexample and explain which system contract or ownership invariant it violates.

\subsection{Practice and evidence}
Compare interrupted blocking I/O with and without SA\_RESTART. Record an expected result, a negative test, observed errno or pthread error code as applicable, and a cleanup trace.

\textbf{Review prompts.} Which state is process-local versus shared? What changes after an interruption, partial completion or concurrent operation?

\section{Chapter 22: Signals: Advanced Features}
\textbf{Textbook entry:} printed page 447.\quad \textbf{Study label:} 2010 textbook foundation.

\subsection{Conceptual map}
\begin{itemize}
\item Realtime signals.
\item Synchronous signal wait.
\item Signalfd.
\end{itemize}
\subsection{API and implementation audit}
Trace the purpose, argument lifetime, return-value convention and failure path for: sigwaitinfo, sigsuspend, signalfd. Identify whether each name represents a kernel system call, a libc wrapper or a library-level operation on your platform; do not infer the category from naming alone.

\subsection{Failure mode and reasoning}
Racing between checking state and sleeping for a signal. Construct a counterexample and explain which system contract or ownership invariant it violates.

\subsection{Practice and evidence}
Implement a sigwait thread with signals blocked across worker threads. Record an expected result, a negative test, observed errno or pthread error code as applicable, and a cleanup trace.

\textbf{Review prompts.} Which state is process-local versus shared? What changes after an interruption, partial completion or concurrent operation?

\section{Chapter 23: Timers and Sleeping}
\textbf{Textbook entry:} printed page 479.\quad \textbf{Study label:} 2010 textbook foundation.

\subsection{Conceptual map}
\begin{itemize}
\item Relative/absolute timeouts.
\item High-resolution sleeps.
\item Timerfd.
\end{itemize}
\subsection{API and implementation audit}
Trace the purpose, argument lifetime, return-value convention and failure path for: clock\_nanosleep, timer\_create, timerfd\_create. Identify whether each name represents a kernel system call, a libc wrapper or a library-level operation on your platform; do not infer the category from naming alone.

\subsection{Failure mode and reasoning}
Building periodic deadlines by accumulating drift. Construct a counterexample and explain which system contract or ownership invariant it violates.

\subsection{Practice and evidence}
Measure periodic deadline error on CLOCK\_MONOTONIC. Record an expected result, a negative test, observed errno or pthread error code as applicable, and a cleanup trace.

\textbf{Review prompts.} Which state is process-local versus shared? What changes after an interruption, partial completion or concurrent operation?

\chapter{Processes and pthreads}
\section{Chapter 24: Process Creation}
\textbf{Textbook entry:} printed page 513.\quad \textbf{Study label:} 2010 textbook foundation.

\subsection{Conceptual map}
\begin{itemize}
\item Fork semantics.
\item Copy-on-write.
\item Descriptor inheritance.
\end{itemize}
\subsection{API and implementation audit}
Trace the purpose, argument lifetime, return-value convention and failure path for: fork, vfork, pipe2. Identify whether each name represents a kernel system call, a libc wrapper or a library-level operation on your platform; do not infer the category from naming alone.

\subsection{Failure mode and reasoning}
Assuming forked memory modifications are shared by default. Construct a counterexample and explain which system contract or ownership invariant it violates.

\subsection{Practice and evidence}
Build a parent-child synchronization pipe and verify copy-on-write. Record an expected result, a negative test, observed errno or pthread error code as applicable, and a cleanup trace.

\textbf{Review prompts.} Which state is process-local versus shared? What changes after an interruption, partial completion or concurrent operation?

\section{Chapter 25: Process Termination}
\textbf{Textbook entry:} printed page 531.\quad \textbf{Study label:} 2010 textbook foundation.

\subsection{Conceptual map}
\begin{itemize}
\item Exit versus \_exit.
\item Handlers.
\item Buffered stdio after fork.
\end{itemize}
\subsection{API and implementation audit}
Trace the purpose, argument lifetime, return-value convention and failure path for: exit, \_exit, atexit. Identify whether each name represents a kernel system call, a libc wrapper or a library-level operation on your platform; do not infer the category from naming alone.

\subsection{Failure mode and reasoning}
Calling exit in a fork child before exec and flushing inherited buffers. Construct a counterexample and explain which system contract or ownership invariant it violates.

\subsection{Practice and evidence}
Reproduce duplicate buffered output and correct it with \_exit. Record an expected result, a negative test, observed errno or pthread error code as applicable, and a cleanup trace.

\textbf{Review prompts.} Which state is process-local versus shared? What changes after an interruption, partial completion or concurrent operation?

\section{Chapter 26: Monitoring Child Processes}
\textbf{Textbook entry:} printed page 541.\quad \textbf{Study label:} 2010 textbook foundation.

\subsection{Conceptual map}
\begin{itemize}
\item Waitpid.
\item Waitid.
\item Sigchld.
\item Zombies and orphans.
\end{itemize}
\subsection{API and implementation audit}
Trace the purpose, argument lifetime, return-value convention and failure path for: waitpid, waitid, sigaction. Identify whether each name represents a kernel system call, a libc wrapper or a library-level operation on your platform; do not infer the category from naming alone.

\subsection{Failure mode and reasoning}
Failing to reap children and leaking zombie process entries. Construct a counterexample and explain which system contract or ownership invariant it violates.

\subsection{Practice and evidence}
Create multiple children and implement a deterministic reaping loop. Record an expected result, a negative test, observed errno or pthread error code as applicable, and a cleanup trace.

\textbf{Review prompts.} Which state is process-local versus shared? What changes after an interruption, partial completion or concurrent operation?

\section{Chapter 27: Program Execution}
\textbf{Textbook entry:} printed page 563.\quad \textbf{Study label:} 2010 textbook foundation.

\subsection{Conceptual map}
\begin{itemize}
\item Execve environment.
\item Search paths.
\item Close-on-exec.
\item Safe launch.
\end{itemize}
\subsection{API and implementation audit}
Trace the purpose, argument lifetime, return-value convention and failure path for: execve, fexecve, posix\_spawnp. Identify whether each name represents a kernel system call, a libc wrapper or a library-level operation on your platform; do not infer the category from naming alone.

\subsection{Failure mode and reasoning}
Constructing an executable command through a shell with untrusted input. Construct a counterexample and explain which system contract or ownership invariant it violates.

\subsection{Practice and evidence}
Launch a fixed binary with a controlled argv and environment. Record an expected result, a negative test, observed errno or pthread error code as applicable, and a cleanup trace.

\textbf{Review prompts.} Which state is process-local versus shared? What changes after an interruption, partial completion or concurrent operation?

\section{Chapter 28: Process Creation and Program Execution in More Detail}
\textbf{Textbook entry:} printed page 591.\quad \textbf{Study label:} 2010 textbook foundation.

\subsection{Conceptual map}
\begin{itemize}
\item Clone.
\item Process attributes across exec/fork.
\item Accounting.
\end{itemize}
\subsection{API and implementation audit}
Trace the purpose, argument lifetime, return-value convention and failure path for: clone, clone3, waitpid. Identify whether each name represents a kernel system call, a libc wrapper or a library-level operation on your platform; do not infer the category from naming alone.

\subsection{Failure mode and reasoning}
Assuming Linux clone flags have POSIX equivalents. Construct a counterexample and explain which system contract or ownership invariant it violates.

\subsection{Practice and evidence}
Record a compatibility matrix and test optional clone3 availability. Record an expected result, a negative test, observed errno or pthread error code as applicable, and a cleanup trace.

\textbf{Review prompts.} Which state is process-local versus shared? What changes after an interruption, partial completion or concurrent operation?

\section{Chapter 29: Threads: Introduction}
\textbf{Textbook entry:} printed page 617.\quad \textbf{Study label:} 2010 textbook foundation.

\subsection{Conceptual map}
\begin{itemize}
\item Pthread lifecycle.
\item Joinability.
\item Shared memory.
\item Thread ids.
\end{itemize}
\subsection{API and implementation audit}
Trace the purpose, argument lifetime, return-value convention and failure path for: pthread\_create, pthread\_join, pthread\_detach. Identify whether each name represents a kernel system call, a libc wrapper or a library-level operation on your platform; do not infer the category from naming alone.

\subsection{Failure mode and reasoning}
Losing ownership of work buffers passed to new threads. Construct a counterexample and explain which system contract or ownership invariant it violates.

\subsection{Practice and evidence}
Run a pool of workers and enforce join and cleanup invariants. Record an expected result, a negative test, observed errno or pthread error code as applicable, and a cleanup trace.

\textbf{Review prompts.} Which state is process-local versus shared? What changes after an interruption, partial completion or concurrent operation?

\section{Chapter 30: Threads: Thread Synchronization}
\textbf{Textbook entry:} printed page 631.\quad \textbf{Study label:} 2010 textbook foundation.

\subsection{Conceptual map}
\begin{itemize}
\item Mutexes.
\item Condition predicates.
\item Wait loops.
\item Deadlock.
\end{itemize}
\subsection{API and implementation audit}
Trace the purpose, argument lifetime, return-value convention and failure path for: pthread\_mutex\_lock, pthread\_cond\_wait, pthread\_cond\_signal. Identify whether each name represents a kernel system call, a libc wrapper or a library-level operation on your platform; do not infer the category from naming alone.

\subsection{Failure mode and reasoning}
Using if instead of while around condition-variable predicates. Construct a counterexample and explain which system contract or ownership invariant it violates.

\subsection{Practice and evidence}
Implement a bounded queue with spurious wakeup-resistant waits. Record an expected result, a negative test, observed errno or pthread error code as applicable, and a cleanup trace.

\textbf{Review prompts.} Which state is process-local versus shared? What changes after an interruption, partial completion or concurrent operation?

\section{Chapter 31: Threads: Thread Safety and Per-Thread Storage}
\textbf{Textbook entry:} printed page 655.\quad \textbf{Study label:} 2010 textbook foundation.

\subsection{Conceptual map}
\begin{itemize}
\item Reentrancy.
\item Once initialization.
\item Thread-specific data.
\end{itemize}
\subsection{API and implementation audit}
Trace the purpose, argument lifetime, return-value convention and failure path for: pthread\_once, pthread\_key\_create, pthread\_getspecific. Identify whether each name represents a kernel system call, a libc wrapper or a library-level operation on your platform; do not infer the category from naming alone.

\subsection{Failure mode and reasoning}
Returning pointers to shared static scratch buffers. Construct a counterexample and explain which system contract or ownership invariant it violates.

\subsection{Practice and evidence}
Make a library function thread-safe with caller-owned storage. Record an expected result, a negative test, observed errno or pthread error code as applicable, and a cleanup trace.

\textbf{Review prompts.} Which state is process-local versus shared? What changes after an interruption, partial completion or concurrent operation?

\section{Chapter 32: Threads: Thread Cancellation}
\textbf{Textbook entry:} printed page 671.\quad \textbf{Study label:} 2010 textbook foundation.

\subsection{Conceptual map}
\begin{itemize}
\item Cancellation states.
\item Cleanup handlers.
\item Cancellation points.
\end{itemize}
\subsection{API and implementation audit}
Trace the purpose, argument lifetime, return-value convention and failure path for: pthread\_cancel, pthread\_cleanup\_push, pthread\_testcancel. Identify whether each name represents a kernel system call, a libc wrapper or a library-level operation on your platform; do not infer the category from naming alone.

\subsection{Failure mode and reasoning}
Leaving mutexes locked after a cancellation point. Construct a counterexample and explain which system contract or ownership invariant it violates.

\subsection{Practice and evidence}
Add deferred cancellation and a cleanup handler to a worker. Record an expected result, a negative test, observed errno or pthread error code as applicable, and a cleanup trace.

\textbf{Review prompts.} Which state is process-local versus shared? What changes after an interruption, partial completion or concurrent operation?

\section{Chapter 33: Threads: Further Details}
\textbf{Textbook entry:} printed page 681.\quad \textbf{Study label:} 2010 textbook foundation.

\subsection{Conceptual map}
\begin{itemize}
\item Thread signal masks.
\item Process-wide behavior.
\item Linux nptl.
\end{itemize}
\subsection{API and implementation audit}
Trace the purpose, argument lifetime, return-value convention and failure path for: pthread\_sigmask, pthread\_kill, pthread\_attr\_setstacksize. Identify whether each name represents a kernel system call, a libc wrapper or a library-level operation on your platform; do not infer the category from naming alone.

\subsection{Failure mode and reasoning}
Delivering application control signals to arbitrary worker threads. Construct a counterexample and explain which system contract or ownership invariant it violates.

\subsection{Practice and evidence}
Centralize signal handling in a dedicated sigwait thread. Record an expected result, a negative test, observed errno or pthread error code as applicable, and a cleanup trace.

\textbf{Review prompts.} Which state is process-local versus shared? What changes after an interruption, partial completion or concurrent operation?

\chapter{Services, credentials and linking}
\section{Chapter 34: Process Groups, Sessions, and Job Control}
\textbf{Textbook entry:} printed page 699.\quad \textbf{Study label:} 2010 textbook foundation.

\subsection{Conceptual map}
\begin{itemize}
\item Sessions.
\item Process groups.
\item Terminal foreground ownership.
\end{itemize}
\subsection{API and implementation audit}
Trace the purpose, argument lifetime, return-value convention and failure path for: setsid, setpgid, tcsetpgrp. Identify whether each name represents a kernel system call, a libc wrapper or a library-level operation on your platform; do not infer the category from naming alone.

\subsection{Failure mode and reasoning}
Treating a process group identifier as a security boundary. Construct a counterexample and explain which system contract or ownership invariant it violates.

\subsection{Practice and evidence}
Build a test shell fragment with jobs and observe group IDs. Record an expected result, a negative test, observed errno or pthread error code as applicable, and a cleanup trace.

\textbf{Review prompts.} Which state is process-local versus shared? What changes after an interruption, partial completion or concurrent operation?

\section{Chapter 35: Process Priorities and Scheduling}
\textbf{Textbook entry:} printed page 733.\quad \textbf{Study label:} 2010 textbook foundation.

\subsection{Conceptual map}
\begin{itemize}
\item Scheduling policy.
\item Nice.
\item Real-time limits.
\item Affinity.
\end{itemize}
\subsection{API and implementation audit}
Trace the purpose, argument lifetime, return-value convention and failure path for: nice, sched\_\allowbreak{}setscheduler, sched\_setaffinity. Identify whether each name represents a kernel system call, a libc wrapper or a library-level operation on your platform; do not infer the category from naming alone.

\subsection{Failure mode and reasoning}
Assuming real-time scheduling is safe or unprivileged. Construct a counterexample and explain which system contract or ownership invariant it violates.

\subsection{Practice and evidence}
Measure scheduler behavior on a non-real-time workload. Record an expected result, a negative test, observed errno or pthread error code as applicable, and a cleanup trace.

\textbf{Review prompts.} Which state is process-local versus shared? What changes after an interruption, partial completion or concurrent operation?

\section{Chapter 36: Process Resources}
\textbf{Textbook entry:} printed page 753.\quad \textbf{Study label:} 2010 textbook foundation.

\subsection{Conceptual map}
\begin{itemize}
\item Resource usage.
\item Rlimit bounds.
\item Cpu/memory control.
\end{itemize}
\subsection{API and implementation audit}
Trace the purpose, argument lifetime, return-value convention and failure path for: getrlimit, setrlimit, getrusage. Identify whether each name represents a kernel system call, a libc wrapper or a library-level operation on your platform; do not infer the category from naming alone.

\subsection{Failure mode and reasoning}
Leaving file-descriptor or address-space growth unlimited. Construct a counterexample and explain which system contract or ownership invariant it violates.

\subsection{Practice and evidence}
Apply RLIMIT\_NOFILE in a child and test predictable failure. Record an expected result, a negative test, observed errno or pthread error code as applicable, and a cleanup trace.

\textbf{Review prompts.} Which state is process-local versus shared? What changes after an interruption, partial completion or concurrent operation?

\section{Chapter 37: Daemons}
\textbf{Textbook entry:} printed page 767.\quad \textbf{Study label:} 2010 textbook foundation.

\subsection{Conceptual map}
\begin{itemize}
\item Daemon lifecycle.
\item Sighup.
\item Logging.
\item Service supervision.
\end{itemize}
\subsection{API and implementation audit}
Trace the purpose, argument lifetime, return-value convention and failure path for: setsid, syslog, openlog. Identify whether each name represents a kernel system call, a libc wrapper or a library-level operation on your platform; do not infer the category from naming alone.

\subsection{Failure mode and reasoning}
Double-forking by default under an external process supervisor. Construct a counterexample and explain which system contract or ownership invariant it violates.

\subsection{Practice and evidence}
Implement a foreground service with explicit shutdown semantics. Record an expected result, a negative test, observed errno or pthread error code as applicable, and a cleanup trace.

\textbf{Review prompts.} Which state is process-local versus shared? What changes after an interruption, partial completion or concurrent operation?

\section{Chapter 38: Writing Secure Privileged Programs}
\textbf{Textbook entry:} printed page 783.\quad \textbf{Study label:} 2010 textbook foundation.

\subsection{Conceptual map}
\begin{itemize}
\item Least privilege.
\item Hostile environment.
\item Toctou.
\item Fail-closed checks.
\end{itemize}
\subsection{API and implementation audit}
Trace the purpose, argument lifetime, return-value convention and failure path for: setresuid, openat, fstat. Identify whether each name represents a kernel system call, a libc wrapper or a library-level operation on your platform; do not infer the category from naming alone.

\subsection{Failure mode and reasoning}
Trusting PATH, cwd, temporary directory, or pathname after checking. Construct a counterexample and explain which system contract or ownership invariant it violates.

\subsection{Practice and evidence}
Write a privileged-operation checklist for a file-processing helper. Record an expected result, a negative test, observed errno or pthread error code as applicable, and a cleanup trace.

\textbf{Review prompts.} Which state is process-local versus shared? What changes after an interruption, partial completion or concurrent operation?

\section{Chapter 39: Capabilities}
\textbf{Textbook entry:} printed page 797.\quad \textbf{Study label:} 2010 textbook foundation.

\subsection{Conceptual map}
\begin{itemize}
\item Capability sets.
\item Exec transitions.
\item Bounding behavior.
\end{itemize}
\subsection{API and implementation audit}
Trace the purpose, argument lifetime, return-value convention and failure path for: capget, capset, prctl. Identify whether each name represents a kernel system call, a libc wrapper or a library-level operation on your platform; do not infer the category from naming alone.

\subsection{Failure mode and reasoning}
Equating a capability with complete containment. Construct a counterexample and explain which system contract or ownership invariant it violates.

\subsection{Practice and evidence}
Inspect capability sets and design a minimal CAP\_NET\_BIND\_SERVICE scenario. Record an expected result, a negative test, observed errno or pthread error code as applicable, and a cleanup trace.

\textbf{Review prompts.} Which state is process-local versus shared? What changes after an interruption, partial completion or concurrent operation?

\section{Chapter 40: Login Accounting}
\textbf{Textbook entry:} printed page 817.\quad \textbf{Study label:} 2010 textbook foundation.

\subsection{Conceptual map}
\begin{itemize}
\item Session records.
\item Utmp and wtmp.
\item Historical login accounting.
\end{itemize}
\subsection{API and implementation audit}
Trace the purpose, argument lifetime, return-value convention and failure path for: getutxent, setutxent, getlogin\_r. Identify whether each name represents a kernel system call, a libc wrapper or a library-level operation on your platform; do not infer the category from naming alone.

\subsection{Failure mode and reasoning}
Assuming utmp accurately captures sessions in modern containers. Construct a counterexample and explain which system contract or ownership invariant it violates.

\subsection{Practice and evidence}
Compare system login records with systemd login sessions, noting variation. Record an expected result, a negative test, observed errno or pthread error code as applicable, and a cleanup trace.

\textbf{Review prompts.} Which state is process-local versus shared? What changes after an interruption, partial completion or concurrent operation?

\section{Chapter 41: Fundamentals of Shared Libraries}
\textbf{Textbook entry:} printed page 833.\quad \textbf{Study label:} 2010 textbook foundation.

\subsection{Conceptual map}
\begin{itemize}
\item Elf dynamic linking.
\item Soname.
\item Rpath/runpath.
\item Pic.
\end{itemize}
\subsection{API and implementation audit}
Trace the purpose, argument lifetime, return-value convention and failure path for: dlopen, readelf, ldd. Identify whether each name represents a kernel system call, a libc wrapper or a library-level operation on your platform; do not infer the category from naming alone.

\subsection{Failure mode and reasoning}
Assuming loaded libraries come from trusted directories. Construct a counterexample and explain which system contract or ownership invariant it violates.

\subsection{Practice and evidence}
Build versioned shared libraries and inspect dynamic sections. Record an expected result, a negative test, observed errno or pthread error code as applicable, and a cleanup trace.

\textbf{Review prompts.} Which state is process-local versus shared? What changes after an interruption, partial completion or concurrent operation?

\section{Chapter 42: Advanced Features of Shared Libraries}
\textbf{Textbook entry:} printed page 859.\quad \textbf{Study label:} 2010 textbook foundation.

\subsection{Conceptual map}
\begin{itemize}
\item Dynamic loading.
\item Symbol visibility.
\item Plugin abi.
\end{itemize}
\subsection{API and implementation audit}
Trace the purpose, argument lifetime, return-value convention and failure path for: dlopen, dlsym, dlclose. Identify whether each name represents a kernel system call, a libc wrapper or a library-level operation on your platform; do not infer the category from naming alone.

\subsection{Failure mode and reasoning}
Loading plugins from attacker-controlled locations. Construct a counterexample and explain which system contract or ownership invariant it violates.

\subsection{Practice and evidence}
Implement a narrowly defined plugin ABI and reject version mismatches. Record an expected result, a negative test, observed errno or pthread error code as applicable, and a cleanup trace.

\textbf{Review prompts.} Which state is process-local versus shared? What changes after an interruption, partial completion or concurrent operation?

\chapter{IPC and memory mapping}
\section{Chapter 43: Interprocess Communication Overview}
\textbf{Textbook entry:} printed page 877.\quad \textbf{Study label:} 2010 textbook foundation.

\subsection{Conceptual map}
\begin{itemize}
\item Ipc choice.
\item Isolation.
\item Throughput, latency and complexity.
\end{itemize}
\subsection{API and implementation audit}
Trace the purpose, argument lifetime, return-value convention and failure path for: pipe, socketpair, mmap, sem\_init. Identify whether each name represents a kernel system call, a libc wrapper or a library-level operation on your platform; do not infer the category from naming alone.

\subsection{Failure mode and reasoning}
Selecting IPC by speed alone without correctness or portability tests. Construct a counterexample and explain which system contract or ownership invariant it violates.

\subsection{Practice and evidence}
Write an IPC architecture decision record for three workloads. Record an expected result, a negative test, observed errno or pthread error code as applicable, and a cleanup trace.

\textbf{Review prompts.} Which state is process-local versus shared? What changes after an interruption, partial completion or concurrent operation?

\section{Chapter 44: Pipes and FIFOs}
\textbf{Textbook entry:} printed page 889.\quad \textbf{Study label:} 2010 textbook foundation.

\subsection{Conceptual map}
\begin{itemize}
\item Pipes, fifos, pipe\_buf atomicity, eof and sigpipe.
\end{itemize}
\subsection{API and implementation audit}
Trace the purpose, argument lifetime, return-value convention and failure path for: pipe2, mkfifo, read, write. Identify whether each name represents a kernel system call, a libc wrapper or a library-level operation on your platform; do not infer the category from naming alone.

\subsection{Failure mode and reasoning}
Forgetting that open write ends keep readers from seeing EOF. Construct a counterexample and explain which system contract or ownership invariant it violates.

\subsection{Practice and evidence}
Build a streaming producer-consumer pipe with deterministic shutdown. Record an expected result, a negative test, observed errno or pthread error code as applicable, and a cleanup trace.

\textbf{Review prompts.} Which state is process-local versus shared? What changes after an interruption, partial completion or concurrent operation?

\section{Chapter 45: Introduction to System V IPC}
\textbf{Textbook entry:} printed page 921.\quad \textbf{Study label:} 2010 textbook foundation.

\subsection{Conceptual map}
\begin{itemize}
\item System v object identity.
\item Ipc namespaces.
\item Lifetime semantics.
\end{itemize}
\subsection{API and implementation audit}
Trace the purpose, argument lifetime, return-value convention and failure path for: ftok, ipcs, ipcrm. Identify whether each name represents a kernel system call, a libc wrapper or a library-level operation on your platform; do not infer the category from naming alone.

\subsection{Failure mode and reasoning}
Leaving persistent IPC objects after program failure. Construct a counterexample and explain which system contract or ownership invariant it violates.

\subsection{Practice and evidence}
Inventory System V IPC objects in a disposable namespace. Record an expected result, a negative test, observed errno or pthread error code as applicable, and a cleanup trace.

\textbf{Review prompts.} Which state is process-local versus shared? What changes after an interruption, partial completion or concurrent operation?

\section{Chapter 46: System V Message Queues}
\textbf{Textbook entry:} printed page 937.\quad \textbf{Study label:} 2010 textbook foundation.

\subsection{Conceptual map}
\begin{itemize}
\item System v messages.
\item Types, sizes, permissions.
\end{itemize}
\subsection{API and implementation audit}
Trace the purpose, argument lifetime, return-value convention and failure path for: msgget, msgsnd, msgrcv, msgctl. Identify whether each name represents a kernel system call, a libc wrapper or a library-level operation on your platform; do not infer the category from naming alone.

\subsection{Failure mode and reasoning}
Assuming message queues are unbounded or ordered globally by priority. Construct a counterexample and explain which system contract or ownership invariant it violates.

\subsection{Practice and evidence}
Model bounded producer pressure and clean up IPC objects. Record an expected result, a negative test, observed errno or pthread error code as applicable, and a cleanup trace.

\textbf{Review prompts.} Which state is process-local versus shared? What changes after an interruption, partial completion or concurrent operation?

\section{Chapter 47: System V Semaphores}
\textbf{Textbook entry:} printed page 965.\quad \textbf{Study label:} 2010 textbook foundation.

\subsection{Conceptual map}
\begin{itemize}
\item System v semaphore sets.
\item Atomic semop.
\item Undo caveats.
\end{itemize}
\subsection{API and implementation audit}
Trace the purpose, argument lifetime, return-value convention and failure path for: semget, semop, semctl. Identify whether each name represents a kernel system call, a libc wrapper or a library-level operation on your platform; do not infer the category from naming alone.

\subsection{Failure mode and reasoning}
Assuming SEM\_UNDO provides full transactional recovery. Construct a counterexample and explain which system contract or ownership invariant it violates.

\subsection{Practice and evidence}
Implement a semaphore-coordinated two-process counter. Record an expected result, a negative test, observed errno or pthread error code as applicable, and a cleanup trace.

\textbf{Review prompts.} Which state is process-local versus shared? What changes after an interruption, partial completion or concurrent operation?

\section{Chapter 48: System V Shared Memory}
\textbf{Textbook entry:} printed page 997.\quad \textbf{Study label:} 2010 textbook foundation.

\subsection{Conceptual map}
\begin{itemize}
\item System v shared segment attachment and teardown.
\end{itemize}
\subsection{API and implementation audit}
Trace the purpose, argument lifetime, return-value convention and failure path for: shmget, shmat, shmdt, shmctl. Identify whether each name represents a kernel system call, a libc wrapper or a library-level operation on your platform; do not infer the category from naming alone.

\subsection{Failure mode and reasoning}
Sharing pointers instead of offsets across mappings. Construct a counterexample and explain which system contract or ownership invariant it violates.

\subsection{Practice and evidence}
Use offsets inside a shared region to avoid address assumptions. Record an expected result, a negative test, observed errno or pthread error code as applicable, and a cleanup trace.

\textbf{Review prompts.} Which state is process-local versus shared? What changes after an interruption, partial completion or concurrent operation?

\section{Chapter 49: Memory Mappings}
\textbf{Textbook entry:} printed page 1017.\quad \textbf{Study label:} 2010 textbook foundation.

\subsection{Conceptual map}
\begin{itemize}
\item Map\_shared versus map\_private.
\item File backing.
\item Msync.
\end{itemize}
\subsection{API and implementation audit}
Trace the purpose, argument lifetime, return-value convention and failure path for: mmap, munmap, msync. Identify whether each name represents a kernel system call, a libc wrapper or a library-level operation on your platform; do not infer the category from naming alone.

\subsection{Failure mode and reasoning}
Accessing beyond a truncated mapping and triggering SIGBUS. Construct a counterexample and explain which system contract or ownership invariant it violates.

\subsection{Practice and evidence}
Map a file and test coherent cross-process reads with synchronization. Record an expected result, a negative test, observed errno or pthread error code as applicable, and a cleanup trace.

\textbf{Review prompts.} Which state is process-local versus shared? What changes after an interruption, partial completion or concurrent operation?

\section{Chapter 50: Virtual Memory Operations}
\textbf{Textbook entry:} printed page 1045.\quad \textbf{Study label:} 2010 textbook foundation.

\subsection{Conceptual map}
\begin{itemize}
\item Mlock/madvise/mincore.
\item Address space operations.
\end{itemize}
\subsection{API and implementation audit}
Trace the purpose, argument lifetime, return-value convention and failure path for: madvise, mlock, mincore. Identify whether each name represents a kernel system call, a libc wrapper or a library-level operation on your platform; do not infer the category from naming alone.

\subsection{Failure mode and reasoning}
Assuming memory residency is identical to persistence. Construct a counterexample and explain which system contract or ownership invariant it violates.

\subsection{Practice and evidence}
Benchmark anonymous mapping access patterns and record caveats. Record an expected result, a negative test, observed errno or pthread error code as applicable, and a cleanup trace.

\textbf{Review prompts.} Which state is process-local versus shared? What changes after an interruption, partial completion or concurrent operation?

\section{Chapter 51: Introduction to POSIX IPC}
\textbf{Textbook entry:} printed page 1057.\quad \textbf{Study label:} 2010 textbook foundation.

\subsection{Conceptual map}
\begin{itemize}
\item Posix ipc naming and lifetime.
\item Api design tradeoffs.
\end{itemize}
\subsection{API and implementation audit}
Trace the purpose, argument lifetime, return-value convention and failure path for: mq\_open, sem\_open, shm\_open. Identify whether each name represents a kernel system call, a libc wrapper or a library-level operation on your platform; do not infer the category from naming alone.

\subsection{Failure mode and reasoning}
Assuming IPC APIs share the same namespace or lifetime rules. Construct a counterexample and explain which system contract or ownership invariant it violates.

\subsection{Practice and evidence}
Prepare a comparison table across POSIX and System V facilities. Record an expected result, a negative test, observed errno or pthread error code as applicable, and a cleanup trace.

\textbf{Review prompts.} Which state is process-local versus shared? What changes after an interruption, partial completion or concurrent operation?

\section{Chapter 52: POSIX Message Queues}
\textbf{Textbook entry:} printed page 1063.\quad \textbf{Study label:} 2010 textbook foundation.

\subsection{Conceptual map}
\begin{itemize}
\item Posix messages.
\item Priorities.
\item Notification.
\item Limits.
\end{itemize}
\subsection{API and implementation audit}
Trace the purpose, argument lifetime, return-value convention and failure path for: mq\_open, mq\_send, mq\_receive, mq\_notify. Identify whether each name represents a kernel system call, a libc wrapper or a library-level operation on your platform; do not infer the category from naming alone.

\subsection{Failure mode and reasoning}
Dropping messages under EAGAIN or losing notification registration. Construct a counterexample and explain which system contract or ownership invariant it violates.

\subsection{Practice and evidence}
Build a bounded mailbox with timed receive and explicit cleanup. Record an expected result, a negative test, observed errno or pthread error code as applicable, and a cleanup trace.

\textbf{Review prompts.} Which state is process-local versus shared? What changes after an interruption, partial completion or concurrent operation?

\section{Chapter 53: POSIX Semaphores}
\textbf{Textbook entry:} printed page 1089.\quad \textbf{Study label:} 2010 textbook foundation.

\subsection{Conceptual map}
\begin{itemize}
\item Unnamed/named semaphores.
\item Post/wait.
\item Ownership distinctions.
\end{itemize}
\subsection{API and implementation audit}
Trace the purpose, argument lifetime, return-value convention and failure path for: sem\_init, sem\_wait, sem\_post, sem\_open. Identify whether each name represents a kernel system call, a libc wrapper or a library-level operation on your platform; do not infer the category from naming alone.

\subsection{Failure mode and reasoning}
Treating a semaphore like a mutex with owner tracking. Construct a counterexample and explain which system contract or ownership invariant it violates.

\subsection{Practice and evidence}
Coordinate a bounded resource pool across pthread workers. Record an expected result, a negative test, observed errno or pthread error code as applicable, and a cleanup trace.

\textbf{Review prompts.} Which state is process-local versus shared? What changes after an interruption, partial completion or concurrent operation?

\section{Chapter 54: POSIX Shared Memory}
\textbf{Textbook entry:} printed page 1107.\quad \textbf{Study label:} 2010 textbook foundation.

\subsection{Conceptual map}
\begin{itemize}
\item Shm\_open and ftruncate.
\item Mapping.
\item Naming and cleanup.
\end{itemize}
\subsection{API and implementation audit}
Trace the purpose, argument lifetime, return-value convention and failure path for: shm\_open, ftruncate, mmap, shm\_unlink. Identify whether each name represents a kernel system call, a libc wrapper or a library-level operation on your platform; do not infer the category from naming alone.

\subsection{Failure mode and reasoning}
Publishing a shared object before initializing its protocol state. Construct a counterexample and explain which system contract or ownership invariant it violates.

\subsection{Practice and evidence}
Create a versioned shared-memory header with initialization marker. Record an expected result, a negative test, observed errno or pthread error code as applicable, and a cleanup trace.

\textbf{Review prompts.} Which state is process-local versus shared? What changes after an interruption, partial completion or concurrent operation?

\section{Chapter 55: File Locking}
\textbf{Textbook entry:} printed page 1117.\quad \textbf{Study label:} 2010 textbook foundation.

\subsection{Conceptual map}
\begin{itemize}
\item Advisory locks.
\item Byte ranges.
\item Lock lifetime and portability.
\end{itemize}
\subsection{API and implementation audit}
Trace the purpose, argument lifetime, return-value convention and failure path for: fcntl, flock, lockf. Identify whether each name represents a kernel system call, a libc wrapper or a library-level operation on your platform; do not infer the category from naming alone.

\subsection{Failure mode and reasoning}
Assuming advisory locks prevent all nonparticipating access. Construct a counterexample and explain which system contract or ownership invariant it violates.

\subsection{Practice and evidence}
Test competing record locks from independent processes. Record an expected result, a negative test, observed errno or pthread error code as applicable, and a cleanup trace.

\textbf{Review prompts.} Which state is process-local versus shared? What changes after an interruption, partial completion or concurrent operation?

\chapter{Sockets and advanced I/O}
\section{Chapter 56: Sockets: Introduction}
\textbf{Textbook entry:} printed page 1149.\quad \textbf{Study label:} 2010 textbook foundation.

\subsection{Conceptual map}
\begin{itemize}
\item Socket addresses.
\item Streams versus datagrams.
\item Endpoint lifecycle.
\end{itemize}
\subsection{API and implementation audit}
Trace the purpose, argument lifetime, return-value convention and failure path for: socket, bind, listen, accept, connect. Identify whether each name represents a kernel system call, a libc wrapper or a library-level operation on your platform; do not infer the category from naming alone.

\subsection{Failure mode and reasoning}
Assuming send and recv preserve application message boundaries. Construct a counterexample and explain which system contract or ownership invariant it violates.

\subsection{Practice and evidence}
Design length-prefixed framing with a maximum frame length. Record an expected result, a negative test, observed errno or pthread error code as applicable, and a cleanup trace.

\textbf{Review prompts.} Which state is process-local versus shared? What changes after an interruption, partial completion or concurrent operation?

\section{Chapter 57: Sockets: UNIX Domain}
\textbf{Textbook entry:} printed page 1165.\quad \textbf{Study label:} 2010 textbook foundation.

\subsection{Conceptual map}
\begin{itemize}
\item Af\_unix pathname/abstract names.
\item Descriptor passing.
\item Peer identity.
\end{itemize}
\subsection{API and implementation audit}
Trace the purpose, argument lifetime, return-value convention and failure path for: socketpair, sendmsg, recvmsg, getsockopt. Identify whether each name represents a kernel system call, a libc wrapper or a library-level operation on your platform; do not infer the category from naming alone.

\subsection{Failure mode and reasoning}
Trusting peer-supplied identity instead of kernel credentials. Construct a counterexample and explain which system contract or ownership invariant it violates.

\subsection{Practice and evidence}
Build a UNIX socket service with authenticated local client acceptance. Record an expected result, a negative test, observed errno or pthread error code as applicable, and a cleanup trace.

\textbf{Review prompts.} Which state is process-local versus shared? What changes after an interruption, partial completion or concurrent operation?

\section{Chapter 58: Sockets: Fundamentals of TCP/IP Networks}
\textbf{Textbook entry:} printed page 1179.\quad \textbf{Study label:} 2010 textbook foundation.

\subsection{Conceptual map}
\begin{itemize}
\item Ipv4/ipv6.
\item Ports.
\item Tcp/udp.
\item Network byte order.
\end{itemize}
\subsection{API and implementation audit}
Trace the purpose, argument lifetime, return-value convention and failure path for: inet\_pton, inet\_ntop, htons, ntohl. Identify whether each name represents a kernel system call, a libc wrapper or a library-level operation on your platform; do not infer the category from naming alone.

\subsection{Failure mode and reasoning}
Interpreting multibyte network fields in host byte order. Construct a counterexample and explain which system contract or ownership invariant it violates.

\subsection{Practice and evidence}
Write IPv4 and IPv6 address conversion round-trip tests. Record an expected result, a negative test, observed errno or pthread error code as applicable, and a cleanup trace.

\textbf{Review prompts.} Which state is process-local versus shared? What changes after an interruption, partial completion or concurrent operation?

\section{Chapter 59: Sockets: Internet Domains}
\textbf{Textbook entry:} printed page 1197.\quad \textbf{Study label:} 2010 textbook foundation.

\subsection{Conceptual map}
\begin{itemize}
\item Getaddrinfo.
\item Connect/bind.
\item Address iteration.
\item Dual stack.
\end{itemize}
\subsection{API and implementation audit}
Trace the purpose, argument lifetime, return-value convention and failure path for: getaddrinfo, socket, connect, freeaddrinfo. Identify whether each name represents a kernel system call, a libc wrapper or a library-level operation on your platform; do not infer the category from naming alone.

\subsection{Failure mode and reasoning}
Hard-coding IPv4 or assuming DNS yields one endpoint. Construct a counterexample and explain which system contract or ownership invariant it violates.

\subsection{Practice and evidence}
Implement address-family-independent outbound connection logic. Record an expected result, a negative test, observed errno or pthread error code as applicable, and a cleanup trace.

\textbf{Review prompts.} Which state is process-local versus shared? What changes after an interruption, partial completion or concurrent operation?

\section{Chapter 60: Sockets: Server Design}
\textbf{Textbook entry:} printed page 1239.\quad \textbf{Study label:} 2010 textbook foundation.

\subsection{Conceptual map}
\begin{itemize}
\item Iterative, fork, thread, and event-driven server designs.
\end{itemize}
\subsection{API and implementation audit}
Trace the purpose, argument lifetime, return-value convention and failure path for: accept4, poll, epoll\_wait. Identify whether each name represents a kernel system call, a libc wrapper or a library-level operation on your platform; do not infer the category from naming alone.

\subsection{Failure mode and reasoning}
Blocking all clients on one slow peer or unbounded worker creation. Construct a counterexample and explain which system contract or ownership invariant it violates.

\subsection{Practice and evidence}
Compare per-connection threads to bounded event-driven service. Record an expected result, a negative test, observed errno or pthread error code as applicable, and a cleanup trace.

\textbf{Review prompts.} Which state is process-local versus shared? What changes after an interruption, partial completion or concurrent operation?

\section{Chapter 61: Sockets: Advanced Topics}
\textbf{Textbook entry:} printed page 1253.\quad \textbf{Study label:} 2010 textbook foundation.

\subsection{Conceptual map}
\begin{itemize}
\item Socket options.
\item Shutdown.
\item Ancillary messages.
\item Multicast.
\end{itemize}
\subsection{API and implementation audit}
Trace the purpose, argument lifetime, return-value convention and failure path for: setsockopt, getsockopt, shutdown, sendmsg. Identify whether each name represents a kernel system call, a libc wrapper or a library-level operation on your platform; do not infer the category from naming alone.

\subsection{Failure mode and reasoning}
Ignoring half-close semantics and EOF in stream protocols. Construct a counterexample and explain which system contract or ownership invariant it violates.

\subsection{Practice and evidence}
Exercise half-close, timeouts, and socket option inspection. Record an expected result, a negative test, observed errno or pthread error code as applicable, and a cleanup trace.

\textbf{Review prompts.} Which state is process-local versus shared? What changes after an interruption, partial completion or concurrent operation?

\section{Chapter 62: Terminals}
\textbf{Textbook entry:} printed page 1289.\quad \textbf{Study label:} 2010 textbook foundation.

\subsection{Conceptual map}
\begin{itemize}
\item Termios.
\item Canonical/raw input.
\item Controlling terminals.
\end{itemize}
\subsection{API and implementation audit}
Trace the purpose, argument lifetime, return-value convention and failure path for: tcgetattr, tcsetattr, cfmakeraw. Identify whether each name represents a kernel system call, a libc wrapper or a library-level operation on your platform; do not infer the category from naming alone.

\subsection{Failure mode and reasoning}
Leaving a terminal in raw mode after abnormal exit. Construct a counterexample and explain which system contract or ownership invariant it violates.

\subsection{Practice and evidence}
Build a line input tool with reliable terminal restoration. Record an expected result, a negative test, observed errno or pthread error code as applicable, and a cleanup trace.

\textbf{Review prompts.} Which state is process-local versus shared? What changes after an interruption, partial completion or concurrent operation?

\section{Chapter 63: Alternative I/O Models}
\textbf{Textbook entry:} printed page 1325.\quad \textbf{Study label:} 2010 textbook foundation.

\subsection{Conceptual map}
\begin{itemize}
\item Select/poll/epoll.
\item Readiness.
\item Edge-triggered drains.
\end{itemize}
\subsection{API and implementation audit}
Trace the purpose, argument lifetime, return-value convention and failure path for: select, poll, epoll\_create1, epoll\_ctl, epoll\_wait. Identify whether each name represents a kernel system call, a libc wrapper or a library-level operation on your platform; do not infer the category from naming alone.

\subsection{Failure mode and reasoning}
Failing to drain nonblocking fds after an EPOLLET readiness event. Construct a counterexample and explain which system contract or ownership invariant it violates.

\subsection{Practice and evidence}
Implement a multi-client readiness loop with backpressure. Record an expected result, a negative test, observed errno or pthread error code as applicable, and a cleanup trace.

\textbf{Review prompts.} Which state is process-local versus shared? What changes after an interruption, partial completion or concurrent operation?

\section{Chapter 64: Pseudoterminals}
\textbf{Textbook entry:} printed page 1375.\quad \textbf{Study label:} 2010 textbook foundation.

\subsection{Conceptual map}
\begin{itemize}
\item Pty master/slave.
\item Controlling terminals.
\item Terminal emulation.
\end{itemize}
\subsection{API and implementation audit}
Trace the purpose, argument lifetime, return-value convention and failure path for: posix\_openpt, grantpt, unlockpt, ptsname\_r. Identify whether each name represents a kernel system call, a libc wrapper or a library-level operation on your platform; do not infer the category from naming alone.

\subsection{Failure mode and reasoning}
Trusting terminal escape sequences or leaving orphaned PTY children. Construct a counterexample and explain which system contract or ownership invariant it violates.

\subsection{Practice and evidence}
Build a pseudo-terminal recorder with escape-sequence risk review. Record an expected result, a negative test, observed errno or pthread error code as applicable, and a cleanup trace.

\textbf{Review prompts.} Which state is process-local versus shared? What changes after an interruption, partial completion or concurrent operation?

\chapter{Cross-chapter synthesis assignments}
\begin{enumerate}
\item Build a secure file converter joining Chapters 4, 5, 15 and 18, with directory-fd resolution and atomic publication.
\item Construct a supervised worker process linking Chapters 20--28, 35 and 36; handle startup failure, timeout, child reaping and shutdown.
\item Build a local IPC broker combining Chapters 29--31, 43--44, 53--54 and 57 with bounded queues and graceful termination.
\item Implement an event-driven daemon connecting Chapters 37--39 and 56--63, with bounded client state and an observability runbook.
\end{enumerate}
\begin{engineeringnote}
Chapter titles and printed opening pages are taken from the uploaded 2010 book. The learning objectives, reviews, examples, and recommended practices are original engineering additions.
\end{engineeringnote}

\chapter*{Source and currency notes}
\addcontentsline{toc}{chapter}{Source and currency notes}
\textbf{Foundational source:} Michael Kerrisk, \emph{The Linux Programming Interface: A Linux and UNIX System Programming Handbook}, No Starch Press, 2010, ISBN 978-1-59327-220-3. The textbook's chapter titles and start pages organize the study guide; explanations, tasks, and code in this suite are independently authored educational material, not reproductions of the original text or its example programs.

\textbf{Current reference baseline:} The Open Group / IEEE, \emph{POSIX.1-2024, Issue 8}, \url{https://pubs.opengroup.org/onlinepubs/9799919799/}; Linux man-pages project, \url{https://man7.org/linux/man-pages/}; Linux kernel documentation, \url{https://docs.kernel.org/}. For Linux-only extensions consult the relevant per-interface manual pages and kernel documentation for the target kernel, architecture and libc; availability can differ materially.

\textbf{Scope:} This is an engineering study and implementation companion. It does not claim to restate a normative standard or guarantee a specific Linux kernel's behavior. System calls, libc wrappers and higher-level library functions are explicitly distinguished where material.
\end{document}
